\section{Численное интегрирование уравнений движения}

Система уравнений~\eqref{eq:newton}–\eqref{eq:kin-q} является системой обыкновенных дифференциальных уравнений вида:
\begin{equation}
  \dot{\mathbf{y}} = \mathbf{f}(t, \mathbf{y}), \qquad \mathbf{y}(t_0) = \mathbf{y}_0,
  \label{eq:ode}
\end{equation}
где $\mathbf{y} = (\mathbf{r}, \mathbf{q}, \mathbf{v}, \boldsymbol{\omega})$ — вектор состояния.
Аналитическое решение этой системы доступно лишь в простейших случаях;
для произвольных сцен применяют методы численного интегрирования.

\subsection{Требования к методу интегрирования}

При выборе метода интегрирования для физического симулятора реального времени необходимо учитывать:
\begin{itemize}
  \item \textbf{Вычислительную стоимость:} каждый шаг должен укладываться в бюджет реального времени;
  \item \textbf{Устойчивость:} метод не должен «взрываться» при больших шагах $h$;
  \item \textbf{Симплектичность:} для консервативных систем желательно сохранять энергию в среднем, а не накапливать её;
  \item \textbf{Совместимость с решателем ограничений:} импульсный решатель (раздел~2.4) вычисляет скоростные поправки, метод должен корректно их применять.
\end{itemize}

\subsection{Явный метод Эйлера}

Простейший метод первого порядка точности:
\begin{equation}
  \mathbf{y}_{n+1} = \mathbf{y}_n + h\,\mathbf{f}(t_n, \mathbf{y}_n).
  \label{eq:explicit-euler}
\end{equation}

Для уравнений движения:
\begin{align}
  \mathbf{v}_{n+1} &= \mathbf{v}_n + h\,\mathbf{M}^{-1}\mathbf{F}_n, \\
  \mathbf{r}_{n+1} &= \mathbf{r}_n + h\,\mathbf{v}_n.
\end{align}

Явный метод Эйлера неустойчив при жёстких системах (stiff systems):
для пружинных связей с высокой жёсткостью $k$ устойчивость требует $h < 2/\sqrt{k/m}$,
что приводит к слишком малым шагам.
Кроме того, он нарушает сохранение энергии: энергия монотонно возрастает, что неприемлемо для длительных симуляций.

\subsection{Симплектический метод Эйлера}

Изменение порядка обновления даёт метод, сохраняющий симплектическую структуру гамильтоновой механики:
\begin{align}
  \mathbf{v}_{n+1} &= \mathbf{v}_n + h\,\mathbf{M}^{-1}\mathbf{F}_n, \label{eq:symplectic-v}\\
  \mathbf{r}_{n+1} &= \mathbf{r}_n + h\,\mathbf{v}_{n+1}. \label{eq:symplectic-r}
\end{align}

Принципиальное отличие: скорость обновляется первой, затем новая скорость используется для обновления положения.
Симплектический метод Эйлера:
\begin{itemize}
  \item сохраняет гамильтоновую структуру системы (не накапливает энергию в среднем);
  \item имеет тот же порядок точности $O(h)$, что и явный метод;
  \item обеспечивает условную устойчивость при шагах, типичных для физических симуляторов ($h \leq 10\text{ мс}$);
  \item является стандартным выбором в игровых и робототехнических симуляторах~\cite{catto2005iterative}.
\end{itemize}

\subsection{Velocity Verlet}

Метод Верле в скоростной форме (velocity Verlet) является симплектическим методом второго порядка точности:
\begin{align}
  \mathbf{r}_{n+1} &= \mathbf{r}_n + h\,\mathbf{v}_n + \tfrac{h^2}{2}\,\mathbf{a}_n, \label{eq:verlet-r}\\
  \mathbf{a}_{n+1} &= \mathbf{M}^{-1}\mathbf{F}(\mathbf{r}_{n+1}), \\
  \mathbf{v}_{n+1} &= \mathbf{v}_n + \tfrac{h}{2}\left(\mathbf{a}_n + \mathbf{a}_{n+1}\right). \label{eq:verlet-v}
\end{align}

Метод требует двух вычислений сил на шаг и хранения ускорения от предыдущего шага.
Он популярен в молекулярной динамике и телесных симуляторах для оффлайн-моделирования,
однако в интерактивных симуляторах уступает симплектическому Эйлеру по стоимости:
при использовании итерационного решателя ограничений пересчёт ускорений после каждой итерации нецелесообразен.

\subsection{Метод Рунге–Кутты 4-го порядка}

РК4 имеет четвёртый порядок точности и подходит для оффлайн-симуляций с высокими требованиями к точности:
\begin{align}
  \mathbf{k}_1 &= \mathbf{f}(t_n,\; \mathbf{y}_n), \notag\\
  \mathbf{k}_2 &= \mathbf{f}\!\left(t_n + \tfrac{h}{2},\; \mathbf{y}_n + \tfrac{h}{2}\mathbf{k}_1\right), \notag\\
  \mathbf{k}_3 &= \mathbf{f}\!\left(t_n + \tfrac{h}{2},\; \mathbf{y}_n + \tfrac{h}{2}\mathbf{k}_2\right), \notag\\
  \mathbf{k}_4 &= \mathbf{f}(t_n + h,\; \mathbf{y}_n + h\,\mathbf{k}_3), \notag\\
  \mathbf{y}_{n+1} &= \mathbf{y}_n + \tfrac{h}{6}\left(\mathbf{k}_1 + 2\mathbf{k}_2 + 2\mathbf{k}_3 + \mathbf{k}_4\right).
\end{align}

Метод требует четырёх вычислений правой части на шаг.
В контексте решателя ограничений применение РК4 нетривиально:
контактные силы разрывны, и промежуточные вычисления $\mathbf{k}_2$, $\mathbf{k}_3$ могут давать некорректные конфигурации столкновений.

\subsection{Сравнение и обоснование выбора}

В таблице~\ref{tab:integrators} приведено сравнение методов по ключевым критериям.

\begin{table}[H]
  \centering
  \caption{Сравнение методов численного интегрирования}
  \label{tab:integrators}
  \begin{tabular}{lcccc}
    \toprule
    Метод & Порядок & Симпл. & Выч./шаг & Совм. с PGS \\
    \midrule
    Явный Эйлер      & 1 & Нет & 1 & Да \\
    Симплектич. Эйлер & 1 & Да  & 1 & Да \\
    Velocity Verlet   & 2 & Да  & 2 & Ограниченно \\
    РК4              & 4 & Нет & 4 & Нет \\
    \bottomrule
  \end{tabular}
\end{table}

Для интерактивной симуляции в реальном времени выбран \textbf{симплектический метод Эйлера}.
Он обеспечивает оптимальное соотношение стоимости и качества:
один вызов функции сил на шаг, естественная интеграция с импульсным решателем ограничений,
сохранение симплектической структуры системы.
Порядок точности $O(h)$ достаточен при типичных шагах $h = 1\text{–}2\text{ мс}$.

\subsection{Интегрирование ориентации}

Для кватерниона уравнение~\eqref{eq:kin-q} интегрируется аналогично.
На практике используют первый порядок:
\begin{equation}
  \mathbf{q}_{n+1} = \mathbf{q}_n + \tfrac{h}{2}\,\boldsymbol{\Omega}(\boldsymbol{\omega}_{n+1})\,\mathbf{q}_n,
  \label{eq:quat-integrate}
\end{equation}
после чего кватернион нормализуется: $\mathbf{q}_{n+1} \leftarrow \mathbf{q}_{n+1} / \|\mathbf{q}_{n+1}\|$.

Нормализация компенсирует численный дрейф нормы и является дешёвой операцией ($O(1)$).

\subsection{Выводы по разделу}

Симплектический метод Эйлера выбран как основной метод интегрирования для разрабатываемого физического движка.
Он удовлетворяет всем ключевым требованиям: минимальная вычислительная стоимость, симплектичность, совместимость с импульсным решателем ограничений.
Шаг $h$ является настраиваемым параметром движка с рекомендованным значением 1–2~мс для задач робототехники.
